\section{The Table of Precedence, Occurrence, and Concurrence}
\label{sec:precedence}

\subsection{The table itself}

\rn{91} introduces the table with the exclusivity clause already quoted, and then prints
twenty-eight numbered lines grouped under the four classes. It is the operative core of the
1960 calendar, and this reference gives it complete.

\begingroup\footnotesize
\renewcommand{\arraystretch}{1.1}
\begin{longtable}{@{}r >{\raggedright\arraybackslash}p{0.86\linewidth}@{}}
\toprule
\multicolumn{2}{@{}l}{\textbf{\latin{Tabella dierum liturgicorum secundum ordinem praecedentiae disposita} (\rn{91})}}\\
\midrule
\endfirsthead
\toprule
\multicolumn{2}{@{}l}{\textit{Table of precedence, continued}}\\
\midrule
\endhead
\bottomrule
\endfoot
\multicolumn{2}{@{}l}{\textit{Liturgical days of the I class}}\\
1 & The feast of the Nativity of the Lord, the Sunday of the Resurrection, and the Sunday of Pentecost (I class with an octave). \\
2 & The Sacred Triduum. \\
3 & The feasts of the Epiphany and of the Ascension of the Lord, of the Most Holy Trinity, of Corpus Christi, of the Heart of Jesus, and of Christ the King. \\
4 & The feasts of the Immaculate Conception and of the Assumption of the Blessed Virgin Mary. \\
5 & The vigil and the octave day of the Nativity of the Lord. \\
6 & The Sundays of Advent, of Lent, and of the Passion, and Low Sunday. \\
7 & Ferias of the I class not named above, namely Ash Wednesday and the Monday, Tuesday, and Wednesday of Holy Week. \\
8 & The Commemoration of All the Faithful Departed, which however yields place to an occurring Sunday. \\
9 & The vigil of Pentecost. \\
10 & The days within the octaves of Easter and of Pentecost. \\
11 & Feasts of the I class of the universal Church not named above. \\
12 & Proper feasts of the I class, namely: (1) the feast of the duly constituted principal Patron of \textit{a}) a nation, \textit{b}) a region or province, whether ecclesiastical or civil, \textit{c}) a diocese; (2) the anniversary of the dedication of the cathedral; (3) the feast of the duly constituted principal Patron of a place, town, or city; (4) the feast and anniversary of the dedication of one's own church, or of a public or semi-public oratory taking the place of a church; (5) the Title of one's own church; (6) the feast of the Title of an Order or Congregation; (7) the feast of the canonised Founder of an Order or Congregation; (8) the feast of the duly constituted principal Patron of an Order or Congregation, and of a religious province. \\
13 & Indult feasts of the I class, movable first, then fixed. \\
\midrule
\multicolumn{2}{@{}l}{\textit{Liturgical days of the II class}}\\
14 & Feasts of the Lord of the II class, movable first, then fixed. \\
15 & Sundays of the II class. \\
16 & Feasts of the II class of the universal Church that are not of the Lord. \\
17 & The days within the octave of the Nativity of the Lord. \\
18 & Ferias of the II class, namely: of Advent from 17 to 23 December inclusive, and the Ember ferias of Advent, of Lent, and of September. \\
19 & Proper feasts of the II class, namely: (1) the feast of the duly constituted secondary Patron of \textit{a}) a nation, \textit{b}) a region or province, \textit{c}) a diocese, \textit{d}) a place, town, or city; (2) feasts of Saints or Blessed under \rn{43}\,d; (3) feasts of Saints proper to a church under \rn{45}\,c; (4) the feast of the beatified Founder of an Order or Congregation (\rn{46}\,b); (5) the feast of the duly constituted secondary Patron of an Order or Congregation and of a religious province (\rn{46}\,d); (6) feasts of Saints or Blessed under \rn{46}\,e. \\
20 & Indult feasts of the II class, movable first, then fixed. \\
21 & Vigils of the II class. \\
\midrule
\multicolumn{2}{@{}l}{\textit{Liturgical days of the III class}}\\
22 & Ferias of Lent and of the Passion, from the Thursday after Ash Wednesday to the Saturday before the Second Sunday of the Passion inclusive, the Ember ferias excepted. \\
23 & Feasts of the III class inscribed in particular calendars, and first proper feasts, namely: (1) feasts of Saints or Blessed under \rn{43}\,d; (2) feasts of Blessed proper to a church (\rn{45}\,d); (3) feasts of Saints or Blessed under \rn{46}\,e; then indult feasts, movable first, then fixed. \\
24 & Feasts of the III class inscribed in the calendar of the universal Church, movable first, then fixed. \\
25 & Ferias of Advent to 16 December inclusive, the Ember ferias excepted. \\
26 & The vigil of the III class. \\
\midrule
\multicolumn{2}{@{}l}{\textit{Liturgical days of the IV class}}\\
27 & The Office of Holy Mary on Saturday. \\
28 & Ferias of the IV class. \\
\end{longtable}
\endgroup

Several features of the table are worth naming because they are invisible to anyone who reasons
from class labels alone.

\begin{itemize}[leftmargin=1.3em, itemsep=0.3em]
\item \textbf{Line 8 carries its own defeat.} The Commemoration of All the Faithful Departed
stands above the vigil of Pentecost and above the days within the paschal octaves, yet the line
itself states that it yields to an occurring Sunday — an exception written into the table rather
than left to \rn{16}\,b.
\item \textbf{Feasts of the Lord of the II class outrank Sundays of the II class} (line 14 above
line 15), while other universal II class feasts fall below them (line 16). This is the table's
implementation of \rn{16}\,a, and it is why the substitution rule exists.
\item \textbf{The days within the octave of the Nativity (line 17) outrank the II class ferias
of late Advent (line 18)} — but those ferias have already ended on 23 December, so the ordering
is operative only against Ember ferias of Lent and September, which cannot coincide with the
Christmas octave either. The line is a formal completion rather than a live conflict.
\item \textbf{Lenten ferias outrank every III class feast, universal or particular} (line 22
above lines 23--24), while Advent ferias fall below both (line 25).
\item \textbf{Particular III class feasts outrank universal ones} (line 23 above line 24). This
inverts the usual expectation and is deliberate: a diocese's own Saint takes its own day against
a universal III class feast of the same class.
\item \textbf{Within lines 13, 14, 20, 23, and 24, movable feasts precede fixed ones.} This
tie-break is stated five times and nowhere else.
\end{itemize}

\subsection{Occurrence}

\rn{92} defines occurrence as the meeting of two or more Offices on one and the same day, and
divides it: \emph{accidental} when a movable liturgical day and a fixed liturgical day meet only
at certain intervals of years, \emph{perpetual} when two liturgical days meet every year.

\rn{93} states the effect and enumerates the four possible outcomes: the Office of the day of
lower grade yields to that of higher grade, and this can happen \latin{per minus nobilis
omissionem, aut commemorationem, aut translationem, aut repositionem} — by omission of the less
noble, or by commemoration, or by translation, or by reposition. Which of the four applies is
settled by \rnn{95--102}.

\rn{94} adds a rule that is easy to violate: a commemoration fixed to a day is \emph{not}
transferred or reposited along with a feast that is transferred or reposited; it is made on its
own day or omitted, according to the rubrics.

\subsubsection{Accidental occurrence and translation}

\rn{95} restricts the remedy of translation sharply: \latin{Ius translationis in alium diem ob
occurrentiam accidentalem \dots\ competit solummodo festis I classis} — the right of translation
to another day on account of accidental occurrence belongs solely to feasts of the I class.
Other feasts accidentally impeded by an Office of higher grade are either commemorated or, that
year, entirely omitted, according to the rubrics. And if two feasts of the same divine Person, or
two feasts of the same Saint or Blessed, occur together, the Office is of the one holding the
higher place in the table and the other is omitted.

\rn{96} gives the destination: a I class feast impeded by a day holding a higher place is
transferred to the next following day that is not of the I or II class. Two named exceptions
follow, and both are described as going \latin{tamquam in sedem propriam}, as though to a proper
seat: the Annunciation, when it must be transferred after Easter, goes to the Monday after Low
Sunday; and the Commemoration of All the Faithful Departed, when it occurs on a Sunday, goes to
the following Monday.

\rnn{97--99} handle multiplicity and identity. If several I class feasts occur on the same day,
the one holding the higher place in the table is celebrated and the others are transferred
according to the order in which they stand in the table (\rn{97}). If several I class feasts fall
to be transferred on successive days, the same order is kept, and \latin{in paritate autem
Officium prius impeditum praecedit} — in a tie, the Office impeded first takes precedence
(\rn{98}). And \rn{99}: transferred feasts are of the same grade as in their proper seat.

\subsubsection{Perpetual occurrence and reposition}

\rn{100} grants the right of reposition — the permanent relocation of a feast whose proper day
is impeded every year — more widely than translation, but not universally:

\begin{studybox}{\rn{100}: who may be reposited}
The right of reposition belongs to \emph{all} feasts of the I and II class, and to
\emph{particular} feasts of the III class occurring outside Advent and Lent that are impeded in
the whole diocese, or in the whole Order or Congregation, or in one's own church.\\[0.4em]
By contrast, feasts of the III class of the universal Church that are perpetually impeded in
some particular calendar, and III class feasts of a diocese or of an Order or Congregation that
are perpetually impeded in only some churches, are perpetually either commemorated or entirely
omitted, according to the rubrics.
\end{studybox}

\rn{101} gives the destination in terms that parallel \rn{61}\,c: feasts to be reposited, if of
the I or II class, are assigned to the nearest following day that is not of the I or II class;
if of the III class, to the nearest following day free from other Offices of equal or higher
grade. \rn{102} then makes the relocation permanent in the strong sense: the day to which a
perpetually impeded feast has been reposited \latin{habetur tamquam dies proprius}, is held as
its proper day, on which the feast is celebrated at the same grade as in its proper seat.

The difference between translation and reposition is therefore not merely one of duration.
A translated feast is displaced for one year and returns; a reposited feast acquires a new
proper day. This is why \rn{61}\,a can speak of a feast being celebrated on a day other than the
\latin{dies natalicius} without any sense of irregularity.

\subsection{Concurrence}

\rn{103} defines concurrence as the meeting of the Vespers of the current liturgical day with
the First Vespers of the following one. The rule is two sentences.

\rn{104}: in concurrence the Vespers of the liturgical day of the higher class are preferred,
and the others are commemorated or not, according to the rubrics. \rn{105}: when the liturgical
days whose Vespers concur are of the same class, \latin{dicuntur integrae secundae Vesperae de
Officio currenti et fit commemoratio sequentis} — the entire Second Vespers of the current
Office are said and a commemoration of the following is made.

Two consequences follow that a Missal-only reader will miss. First, concurrence is decided by
\emph{class}, not by position in the table of precedence: \rn{104} says \latin{classis
superioris}, and only when the classes are equal does \rn{105} settle the matter in favour of
the day in possession. This is a genuinely different rule from occurrence, where \rn{7} makes the
table exclusive. Second, only days that have First Vespers can concur at all — I class feasts
(\rn{37}\,a), Sundays (\rn{13}), and II class feasts of the Lord that have acquired First Vespers
by taking a Sunday's place (\rn{37}\,c).

The Acta print, after the code, a \latin{Tabella occurrentiae} and a \latin{Tabella
concurrentiae} with three \latin{Notanda}, which restate that a feast of the Lord of the I or II
class occurring on a Sunday takes that Sunday's place with all its rights and privileges; that
of two feasts of the same divine Person or the same Saint the higher in the table is kept and
the other omitted; and that when a feast of the Lord of the I or II class concurs with any
Sunday, or the reverse, the Vespers are ordered by the table of concurrence, but no
commemoration of the concurring Sunday is ever made at the Vespers of the feast of the Lord, nor
conversely. Those tables belong to the Office and are printed in the Breviary, not in the Missal;
a reader working only from the 1962 Missal will not find them, and this reference does not
reproduce them.
